% Rozdział 5 — szkic (wygenerowany 2026-10-17 na podstawie analysis/, commit kampanii 2287a55).
% Wymaga w preambule: \usepackage{booktabs}, \usepackage{graphicx}.
% Tabele i rysunki: skopiuj katalogi analysis/tables/ i analysis/figures/ do Overleafa
% jako tables/ i figures/ (ścieżki poniżej są względne do katalogu głównego projektu).
% Liczby w tekście pochodzą z analysis/results.csv i analysis/cost.csv — po ponownym
% uruchomieniu scripts/analyze.py sprawdź, czy się nie zmieniły.
% Jedyny odsyłacz spoza rozdziału: \ref{ch:metodyka} — dopasuj do etykiety rozdziału 4.

\chapter{Analiza wyników, wykorzystania zasobów i efektywności finansowej}
\label{ch:wyniki}

W rozdziale przedstawiono wyniki kampanii pomiarowej przeprowadzonej według metodyki
opisanej w rozdziale~\ref{ch:metodyka}. Najpierw omówiono przepustowość i~opóźnienia
czterech konfiguracji, następnie wykorzystanie zasobów i~wąskie gardła, wpływ warstwy
dyskowej, fazy przejściowe związane z~mechanizmami kredytowymi oraz zmienność wyników
między wdrożeniami. Rozdział zamyka eksperyment wyjaśniający dotyczący pamięci podręcznej
hosta i~analiza efektywności finansowej, łącząca wydajność z~kosztem utrzymania.

\section{Zakres danych i sposób prezentacji wyników}
\label{sec:wyniki-zakres}

Kampania główna objęła 18 sesji pomiarowych wykonanych między 10 a~17 października 2026~r.
w~regionie Belgium Central: po cztery sesje konfiguracji PaaS General Purpose, IaaS
Premium SSD i~PaaS Burstable oraz sześć sesji konfiguracji IaaS Standard SSD. Każda sesja
była osobnym wdrożeniem środowiska (od \texttt{terraform apply} do \texttt{terraform destroy})
i~obejmowała pięć 10-minutowych przebiegów pomiarowych, co daje 90 przebiegów. Kolejność
sesji w~czterech pierwszych rundach wylosowano przed rozpoczęciem kampanii jako kwadrat
łaciński, a~dodatkowe sesje konfiguracji IaaS Standard SSD wynikają z~reguły zatrzymania
(cel: 95-procentowy przedział ufności średniej przepustowości nie szerszy niż $\pm5\%$
średniej, maksymalnie sześć sesji). Przebiegi pilotażowe, wykonane według wcześniejszej
wersji protokołu, nie wchodzą do analizy.

Jednostką statystyczną jest sesja, a~nie pojedynczy przebieg: przebiegi jednej sesji
korzystają z~tego samego wdrożenia (tych samych maszyn, dysku i~serwera), więc nie są
niezależnymi próbami konfiguracji. Każdą wielkość wyznaczono jako średnią przebiegów
w~sesji, a~następnie średnią tych średnich po sesjach; podawany przedział to 95-procentowy
przedział ufności średniej z~rozkładu t-Studenta o~$k-1$ stopniach swobody, gdzie $k$
oznacza liczbę sesji. Do analizy głównej wchodzą wyłącznie przebiegi spełniające
zdefiniowane przed pilotażem kryterium stanu ustalonego (wyczerpane wszystkie pule
kredytów, które pod tym obciążeniem się wyczerpują). Przebiegi niespełniające kryterium
pozostawiono w~danych i~uwzględniono w~analizie wrażliwości (podrozdział~\ref{sec:wyniki-zmiennosc}).
Pominięto także przebiegi jednej sesji przerwanej awarią po stronie maszyny sterującej
(podrozdział~\ref{sec:wyniki-zmiennosc}); jej miejsce zajęła sesja powtórzona.

\section{Przepustowość i opóźnienia}
\label{sec:wyniki-wydajnosc}

Tabela~\ref{tab:wydajnosc} i~rysunek~\ref{fig:tps} zestawiają przepustowość (liczbę
transakcji na sekundę, TPS) oraz opóźnienia transakcji. Najwyższą przepustowość osiągnęła
konfiguracja PaaS General Purpose: $1358{,}3 \pm 18{,}7$~TPS, o~45\% więcej niż IaaS
Premium SSD ($935{,}5 \pm 43{,}6$~TPS). Konfiguracje ograniczone małym budżetem zasobów
uzyskały wartości o~rząd wielkości niższe: IaaS Standard SSD $172{,}4 \pm 15{,}1$~TPS,
a~PaaS Burstable $121{,}4 \pm 2{,}6$~TPS. Przedziały ufności nie nakładają się dla żadnej
pary konfiguracji, a~najmniejsza różnica (IaaS Standard SSD wobec PaaS Burstable) wynosi
42\%, więc kolejność konfiguracji jest jednoznaczna.

\begin{table}[htbp]
  \centering
  \caption{Przepustowość i opóźnienia transakcji (średnia średnich sesji $\pm$ 95\% przedział ufności; $k$ -- liczba sesji)}
  \label{tab:wydajnosc}
  \small
  \input{tables/performance.tex}
\end{table}

\begin{figure}[htbp]
  \centering
  \includegraphics[width=\textwidth]{figures/tps.pdf}
  \caption{Przepustowość konfiguracji: puste punkty -- średnie poszczególnych sesji, pełny punkt -- średnia z~95\% przedziałem ufności (skala logarytmiczna)}
  \label{fig:tps}
\end{figure}

Średnie opóźnienie transakcji jest w~każdej konfiguracji zgodne z~prawem Little'a dla
stałej liczby 25 klientów ($\bar{L} \approx 25/\mathrm{TPS}$), co potwierdza, że klienci
pracowali w~pętli zamkniętej bez czasu oczekiwania, a~przepustowość i~średnie opóźnienie
niosą tę samą informację. Różnią się natomiast ogony rozkładu (rysunek~\ref{fig:latencja}).
W~konfiguracji PaaS General Purpose 99. percentyl wynosi $47{,}3$~ms, a~w~IaaS Premium SSD
$69{,}0$~ms. W~konfiguracji IaaS Standard SSD p99 sięga ok. 400~ms, a~w~PaaS Burstable
przekracza sekundę ($1012{,}7 \pm 109{,}7$~ms), przy p50 tylko ok. 115--120~ms: rozkład
opóźnień jest tam silnie prawostronnie skośny. Wynika to z~mechanizmu kredytowego CPU
(podrozdział~\ref{sec:wyniki-bursting}): w~stanie wyczerpanych kredytów moc obliczeniowa
jest przydzielana nierównomiernie, co wydłuża czas obsługi części transakcji. Przedziały
ufności dla p99 są szersze niż dla wartości średnich (np. $\pm 13\%$ dla IaaS Premium SSD),
ponieważ percentyl ogonowy jest wrażliwy na pojedyncze zdarzenia w~oknie pomiaru, takie
jak zapis checkpointu.

\begin{figure}[htbp]
  \centering
  \includegraphics[width=\textwidth]{figures/latency.pdf}
  \caption{Percentyle opóźnienia transakcji p50, p95 i~p99 (średnie z~95\% przedziałami ufności, skala logarytmiczna)}
  \label{fig:latencja}
\end{figure}

\section{Wykorzystanie zasobów i wąskie gardła}
\label{sec:wyniki-zasoby}

Tabela~\ref{tab:zasoby} przedstawia wykorzystanie zasobów w~oknach pomiarowych według
metryk Azure Monitor, a~tabela~\ref{tab:latencja-io} -- opóźnienia operacji wejścia/wyjścia
mierzone po stronie platformy i~wewnątrz PostgreSQL (widok \texttt{pg\_stat\_io}).
Zestawienie pozwala wskazać wąskie gardło każdej konfiguracji.

\begin{table}[htbp]
  \centering
  \caption{Wykorzystanie zasobów w~oknie pomiaru: IOPS dysku danych, liczba operacji I/O na transakcję, wykorzystanie limitu IOPS dysku, CPU i~szczytowe wykorzystanie RAM}
  \label{tab:zasoby}
  \resizebox{\textwidth}{!}{\input{tables/resources.tex}}
\end{table}

\begin{table}[htbp]
  \centering
  \caption{Opóźnienia operacji I/O: latencja dysku według Azure Monitor (dostępna tylko dla maszyn wirtualnych), średni czas odczytu i~zapisu bloku według \texttt{pg\_stat\_io} oraz udział trafień w~buforze \texttt{shared\_buffers}}
  \label{tab:latencja-io}
  \resizebox{\textwidth}{!}{\input{tables/io-latency.tex}}
\end{table}

\begin{itemize}
  \item \textbf{IaaS Premium SSD i~IaaS Standard SSD} są ograniczone przez dysk: wykorzystanie
  limitu IOPS dysku wynosi odpowiednio $98{,}2\%$ i~$99{,}0\%$, przy CPU na poziomie $38{,}8\%$
  i~$11{,}2\%$. Dysk pracuje na poziomie bazowym swojego tieru (2300 i~500~IOPS), a~procesor
  ma zapas.
  \item \textbf{PaaS General Purpose} jest ograniczony przez CPU ($83{,}9\%$), a~dysk pracuje
  na $69{,}9\%$ limitu. Usługa dostarcza przy tym więcej operacji odczytu, niż wynika
  z~limitu dysku: wykorzystanie limitu odpowiada w~przybliżeniu samym zapisom, co wskazuje,
  że odczyty są w~dużej części obsługiwane przez pamięć podręczną hosta. Potwierdza to
  średni czas odczytu bloku: $1{,}60$~ms wobec $5{,}09$~ms na IaaS Premium SSD, mimo
  identycznej wielkości i~klasy dysku.
  \item \textbf{PaaS Burstable} jest ograniczony przez CPU ($84{,}9\%$ mocy przyciętej do
  poziomu bazowego po wyczerpaniu kredytów), a~dysk wykorzystuje zaledwie $8{,}3\%$ limitu.
  Jako jedyna konfiguracja ma też wyraźnie dłuższy czas zapisu bloku ($1{,}19$~ms wobec
  setnych części milisekundy w~pozostałych), co jest zgodne z~brakiem mocy obliczeniowej
  na obsługę zapisu.
\end{itemize}

Liczba operacji I/O przypadających na transakcję mieści się w~przedziale 2,5--3,7 i~jest
najwyższa w~konfiguracjach z~najmniejszym budżetem zasobów. Udział trafień w~buforze
\texttt{shared\_buffers} wynosi 92--94\% we wszystkich konfiguracjach, co jest zgodne
z~założeniem projektu: baza o~rozmiarze 14,6~GiB nie mieści się w~8~GiB pamięci serwera,
więc obciążenie wymusza rzeczywiste operacje dyskowe. Wykorzystanie pamięci RAM nie było
ograniczeniem w~żadnej konfiguracji; wartości dla IaaS i~PaaS nie są jednak bezpośrednio
porównywalne, ponieważ pochodzą z~różnych metryk (\texttt{Available Memory Bytes} traktuje
pamięć podręczną systemu plików jako dostępną, \texttt{memory\_percent} usługi zarządzanej
-- jako zajętą), a~serwer PaaS Burstable ma tylko 2~GiB pamięci.

\section{Wpływ warstwy dyskowej}
\label{sec:wyniki-dysk}

Konfiguracje IaaS różnią się wyłącznie tierem dysku danych o~tej samej pojemności 512~GB,
więc ich porównanie izoluje wpływ warstwy dyskowej. Dysk Premium SSD P20 zapewnił
$5{,}4$-krotnie wyższą przepustowość niż Standard SSD E20 ($935{,}5$ wobec $172{,}4$~TPS),
przy $4{,}6$-krotnej różnicy provisionowanych IOPS (2300 wobec 500). Przewaga większa niż
stosunek IOPS wynika z~większej liczby operacji na transakcję na dysku E20 ($3{,}28$ wobec
$2{,}51$): przy małym budżecie IOPS zapis checkpointu zajmuje znaczną część dostępnych
operacji i~ogranicza odczyty potrzebne transakcjom. Efekt ten widać w~przebiegu
przepustowości w~czasie (rysunek~\ref{fig:burnin-e20}) jako oscylacje o~okresie ok. 10~minut,
równym parametrowi \texttt{checkpoint\_timeout}. Okno pomiaru o~długości dokładnie jednego
cyklu checkpointu (600~s) sprawia, że każdy przebieg obejmuje każdą fazę cyklu raz,
niezależnie od momentu startu; kontrola licznikami \texttt{pg\_stat\_bgwriter} potwierdziła
jeden zakończony checkpoint w~89 z~90 okien pomiarowych.

Różnicę widać także w~opóźnieniach operacji dyskowych: latencja dysku według Azure Monitor
wynosi $8{,}66$~ms dla P20 i~$51{,}61$~ms dla E20, a~średni czas odczytu bloku w~PostgreSQL
odpowiednio $5{,}09$ i~$25{,}92$~ms. Ponieważ procesor konfiguracji E20 jest wykorzystany
w~$11\%$, maszyna wirtualna tej samej klasy co w~konfiguracji P20 pozostaje w~dużej mierze
bezczynna -- o~wydajności decyduje wyłącznie tier dysku.

\section{Fazy przejściowe: mechanizmy kredytowe}
\label{sec:wyniki-bursting}

Trzy z~czterech konfiguracji korzystają z~mechanizmów kredytowych, które przez pewien czas
pozwalają przekroczyć poziom bazowy zasobu: pula kredytów burst dysku (IaaS) i~kredyty CPU
(PaaS Burstable). Wynik pomiaru zależy więc od stanu puli przed jego rozpoczęciem. Z~tego
powodu każda sesja zaczynała się adaptacyjnym burn-inem trwającym do wyczerpania pul, które
pod danym obciążeniem się wyczerpują (minimum 60~minut), a~przed każdym przebiegiem
adaptacyjny warm-up ponownie doprowadzał pule do wyczerpania. Tabela~\ref{tab:burnin}
podsumowuje fazy burn-inu wszystkich sesji, a~rysunki~\ref{fig:burnin-p20}--\ref{fig:burnin-b1ms}
pokazują ich przebieg w~pierwszych sesjach kampanii.

\begin{table}[htbp]
  \centering
  \caption{Fazy burn-inu: liczba i~czas trwania burn-inów, moment wyczerpania puli kredytów oraz średnia przepustowość przed wyczerpaniem i~po nim (wartości z~burn-inu, bez resetu stanu bazy -- nie są bezpośrednio porównywalne z~przebiegami pomiarowymi)}
  \label{tab:burnin}
  \small
  \input{tables/burnin.tex}
\end{table}

\begin{figure}[htbp]
  \centering
  \includegraphics[width=\textwidth]{figures/burnin-iaas-premium-ssd.pdf}
  \caption{Burn-in konfiguracji IaaS Premium SSD (pierwsza sesja kampanii): przepustowość w~kolejnych minutach i~zużycie puli kredytów burst dysku}
  \label{fig:burnin-p20}
\end{figure}

W~konfiguracji IaaS Premium SSD faza burstu trwała 32--57 minut i~dawała ok. 1210~TPS;
po wyczerpaniu puli przepustowość spadła do ok. 910~TPS, czyli o~ok. 25\%. Pomiar
wykonany bez burn-inu, bezpośrednio po wdrożeniu, zawyżałby więc wydajność tego tieru
o~jedną trzecią i~opisywałby stan, którego dysk nie jest w~stanie utrzymać. Dysk Standard
SSD E20 ma niewielki margines burstu (600 wobec 500~IOPS), więc spadek po wyczerpaniu puli
był mały (ok. 214 $\to$ 198~TPS, $-7\%$). Wyczerpanie puli następowało zwykle po 31--36
minutach, ale w~jednej sesji dopiero po 145 minutach -- dysk tej instancji wykonywał
mniej operacji ponad poziom bazowy, co adaptacyjny burn-in wykrył i~wydłużył bez
interwencji.

\begin{figure}[htbp]
  \centering
  \includegraphics[width=\textwidth]{figures/burnin-iaas-standard-ssd.pdf}
  \caption{Burn-in konfiguracji IaaS Standard SSD (pierwsza sesja kampanii); widoczne oscylacje przepustowości o~okresie cyklu checkpointu}
  \label{fig:burnin-e20}
\end{figure}

Najdłuższą fazę przejściową ma PaaS Burstable: kredyty CPU zmniejszały się w~tempie
ok. 0,12 na minutę i~wyczerpywały po 196--217 minutach, po czym przepustowość spadała
z~ok. 230 do ok. 150~TPS (rysunek~\ref{fig:burnin-b1ms}). W~pierwszych ok. 30 minutach
widoczna jest dodatkowo krótsza faza burstu operacji I/O po stronie warstwy obliczeniowej,
dla której usługa nie publikuje metryki. Stan utrzymywalny tej konfiguracji zaczyna się
więc dopiero po ponad trzech godzinach ciągłego obciążenia. Konfiguracja PaaS General Purpose
nie ma puli, która by się wyczerpywała; jej burn-in miał stałą długość 60 minut i~stabilną
przepustowość.

\begin{figure}[htbp]
  \centering
  \includegraphics[width=\textwidth]{figures/burnin-paas-burstable.pdf}
  \caption{Burn-in konfiguracji PaaS Burstable (pierwsza sesja kampanii): przepustowość i~pozostałe kredyty CPU}
  \label{fig:burnin-b1ms}
\end{figure}

Znaczenie kryterium stanu ustalonego pokazało zdarzenie z~jednej sesji konfiguracji IaaS
Premium SSD: po wyczerpaniu puli w~trakcie burn-inu metryka zużycia kredytów w~ciągu pięciu
minut spadła ze 100\% do 5\%, czego pod ciągłym obciążeniem nie da się wyjaśnić odnawianiem
puli -- najpewniej nastąpił reset kredytów po stronie platformy. Cztery kolejne przebiegi
mierzyły ponownie dysk w~fazie burstu (do 1391~TPS); kryterium stanu ustalonego oznaczyło
je jako nieważne, a~sesja weszła do analizy głównej jednym przebiegiem.

\section{Zmienność między wdrożeniami i analiza wrażliwości}
\label{sec:wyniki-zmiennosc}

Tabela~\ref{tab:zmiennosc} porównuje odchylenie standardowe średnich sesji (zmienność między
wdrożeniami) z~odchyleniem standardowym przebiegów w~obrębie sesji (powtarzalność pomiaru
na jednym wdrożeniu). W~konfiguracjach PaaS General Purpose i~IaaS Standard SSD zmienność
między wdrożeniami jest większa niż w~obrębie sesji (stosunek $1{,}83$ i~$1{,}36$): kolejne
wdrożenie tej samej konfiguracji trafia na inny sprzęt i~daje inny wynik. W~przypadku IaaS
Standard SSD średnie sesji wahały się od $153{,}6$ do $196{,}8$~TPS, co przy sześciu sesjach
dało przedział ufności $\pm8{,}7\%$ średniej -- jedyny szerszy od założonego celu $\pm5\%$.
Zgodnie z~planem kampanii raportowana jest osiągnięta precyzja. Gdyby jednostką statystyczną
był pojedynczy przebieg, przedziały ufności byłyby zaniżone, ponieważ nie obejmowałyby
zmienności między wdrożeniami. W~konfiguracji PaaS Burstable jest odwrotnie: pojedyncze
przebiegi różnią się znacznie (współczynnik zmienności $8{,}5\%$, efekt krótkich zrywów
przydziału CPU w~stanie wyczerpanych kredytów), ale średnie sesji są bardzo zgodne ($1{,}3\%$).

\begin{table}[htbp]
  \centering
  \caption{Zmienność przepustowości: odchylenie standardowe średnich sesji i~średnie odchylenie standardowe przebiegów w~sesji (TPS), współczynniki zmienności oraz ich stosunek}
  \label{tab:zmiennosc}
  \small
  \input{tables/variance.tex}
\end{table}

Tabela~\ref{tab:wrazliwosc} pokazuje, jak zmieniłaby się średnia przepustowość przy innych
regułach doboru przebiegów. Uwzględnienie przebiegów niespełniających kryterium stanu
ustalonego zmienia wynik tylko dla IaaS Premium SSD (z~$935{,}5$ do $985{,}1$~TPS i~z~przedziałem
$\pm14\%$) -- są to przebiegi zmierzone w~fazie burstu, opisane wyżej. Wykluczenie jedynego
okna bez dokładnie jednego zakończonego checkpointu nie zmienia wyników. Dołączenie
przebiegów sesji przerwanej awarią maszyny sterującej (zapełniony dysk lokalny uniemożliwił
usunięcie środowiska po kompletnym pomiarze) zmienia wynik IaaS Standard SSD w~granicach
przedziału ufności ($169{,}9 \pm 13{,}6$~TPS). Wnioski z~analizy głównej są więc odporne
na przyjęte reguły.

\begin{table}[htbp]
  \centering
  \caption{Analiza wrażliwości średniej przepustowości [TPS] na reguły doboru przebiegów}
  \label{tab:wrazliwosc}
  \resizebox{\textwidth}{!}{\input{tables/sensitivity.tex}}
\end{table}

\section{Eksperyment wyjaśniający: pamięć podręczna hosta}
\label{sec:wyniki-cache}

Wyniki sanity checków wskazywały, że część przewagi PaaS General Purpose nad IaaS Premium SSD
może wynikać z~pamięci podręcznej hosta, z~której korzysta usługa zarządzana, a~której
dysk danych IaaS świadomie nie używał (\texttt{caching = None}). Przed pomiarem sformułowano
hipotezę, że po włączeniu pamięci podręcznej odczytu konfiguracja IaaS Premium SSD zbliży się
do PaaS General Purpose, a~wąskim gardłem stanie się CPU. Eksperyment, wykonany poza główną
macierzą w~jednej ważnej sesji (pierwsza sesja nie dała przebiegów w~stanie ustalonym z~powodu
progu wyczerpania kredytów CPU, skorygowanego na podstawie krzywej kredytów z~burn-inu),
przedstawia tabela~\ref{tab:cache}. Ze względu na jedną sesję wyniki mają charakter opisowy.

\begin{table}[htbp]
  \centering
  \caption{Eksperyment wyjaśniający: IaaS Premium SSD z~pamięcią podręczną odczytu hosta na tle konfiguracji bez niej i~PaaS General Purpose (dla eksperymentu jedna sesja, $n=4$ przebiegi -- bez przedziału ufności)}
  \label{tab:cache}
  \resizebox{\textwidth}{!}{\input{tables/explanatory.tex}}
\end{table}

Pamięć podręczna odczytu zmniejszyła obciążenie dysku z~$98{,}2\%$ do ok. $50\%$ limitu IOPS
i~skróciła średni czas odczytu bloku z~$5{,}09$ do $0{,}96$~ms, a~więc usunęła wąskie gardło
dyskowe. Przepustowość wzrosła jednak tylko o~ok. 3,5\% (do $968{,}4$~TPS), a~wykorzystanie
CPU ($39{,}3\%$) odpowiada poziomowi bazowemu maszyny \texttt{Standard\_B2s\_v2} (40\%),
przy kredytach CPU utrzymujących się na minimalnym poziomie. Nowym wąskim gardłem stał się
więc procesor, ale nie jego pełna moc, tylko poziom bazowy maszyny klasy burstable.
Hipoteza potwierdziła się w~części dotyczącej przesunięcia wąskiego gardła na CPU, nie
potwierdziła się w~części dotyczącej zbliżenia do PaaS General Purpose. Różnica między
modelami w~parze IaaS Premium SSD / PaaS General Purpose wynika więc z~dwóch czynników
naraz: z~pamięci podręcznej hosta po stronie usługi zarządzanej oraz z~klasy jednostki
obliczeniowej (2~vCore bez mechanizmu kredytowego wobec maszyny z~poziomem bazowym 40\%).

\section{Efektywność finansowa}
\label{sec:wyniki-koszty}

Koszt utrzymania wyznaczono z~cennika katalogowego Azure (model pay-as-you-go, region
Belgium Central, USD; migawka cennika z~17 października 2026~r. zapisana w~repozytorium),
uwzględniając wyłącznie zasoby strony bazy danych: maszynę wirtualną albo jednostkę
obliczeniową usługi zarządzanej oraz magazyn danych. Pominięto maszynę kliencką, publiczne
adresy IP służące do administracji i~sieć, ponieważ stanowią narzędzie pomiarowe, a~nie
element konfiguracji. Miesiąc przyjęto jako 730~godzin. Dla dysku Standard SSD uwzględniono
opłatę za transakcje (0,002~USD za 10~tys.), naliczaną od 2023~r. tylko do godzinowego limitu
zależnego od rozmiaru dysku; dla E20 wynosi on 502~tys. transakcji na godzinę. Dysk E20
wykonywał pod tym obciążeniem ok. 560~IOPS (ok. 2~mln operacji na godzinę), więc zawsze
osiągał limit. Rzeczywisty rachunek subskrypcji nie mógł posłużyć za źródło kosztów,
ponieważ w~danych Cost Management nie pojawiła się żadna pozycja usługi Azure Database for
PostgreSQL, a~pozostałe pozycje obejmowały także narzędzie pomiarowe. Składniki kosztu
przedstawia tabela~\ref{tab:koszt-skladniki}.

\begin{table}[htbp]
  \centering
  \caption{Składniki godzinowego kosztu utrzymania strony bazy danych według cennika pay-as-you-go [USD/h]}
  \label{tab:koszt-skladniki}
  \small
  \input{tables/cost-components.tex}
\end{table}

Efektywność finansową wyrażono jako liczbę transakcji przypadających na jednego dolara
kosztu utrzymania, czyli $\mathrm{TPS} \cdot 3600 / c_h$, gdzie $c_h$ to koszt godzinowy;
przedział ufności wynika z~przedziału ufności przepustowości, ponieważ koszt cennikowy jest
stały. Wyniki przedstawiają tabela~\ref{tab:koszt-efektywnosc} i~rysunek~\ref{fig:koszt}.

\begin{table}[htbp]
  \centering
  \caption{Koszt utrzymania i~efektywność finansowa (w~nawiasach 95\% przedziały ufności wynikające z~niepewności przepustowości)}
  \label{tab:koszt-efektywnosc}
  \resizebox{\textwidth}{!}{\input{tables/cost-efficiency.tex}}
\end{table}

\begin{figure}[htbp]
  \centering
  \includegraphics[width=\textwidth]{figures/cost-efficiency.pdf}
  \caption{Liczba transakcji na 1~USD kosztu utrzymania (cennik pay-as-you-go) z~95\% przedziałami ufności}
  \label{fig:koszt}
\end{figure}

Najwyższą efektywność finansową mają konfiguracje PaaS General Purpose ($16{,}35$~mln
transakcji na dolara) i~IaaS Premium SSD ($16{,}07$~mln). Ich przedziały ufności nakładają
się, więc na podstawie pomiarów nie można wskazać, która jest efektywniejsza: PaaS General
Purpose jest o~45\% wydajniejszy i~o~43\% droższy w~utrzymaniu (218 wobec 153~USD miesięcznie).
Usługa zarządzana zapewnia przy tym w~tej cenie funkcje, których konfiguracja IaaS nie
obejmuje i~które wymagałyby dodatkowej pracy administracyjnej (automatyczne kopie zapasowe
z~odtwarzaniem do punktu w~czasie, aktualizacje, monitorowanie). Konfiguracja PaaS Burstable
jest najtańsza w~utrzymaniu (85~USD miesięcznie), ale daje $4{,}3$-krotnie mniej transakcji
na dolara niż PaaS General Purpose; ponad 80\% jej kosztu przypada na magazyn danych
o~pojemności 512~GiB, przyjętej w~projekcie eksperymentu dla zachowania jednakowej warstwy
dyskowej we wszystkich konfiguracjach. Najniższą efektywność finansową ma IaaS Standard SSD
($2{,}46$~mln transakcji na dolara): opłata za transakcje dysku ($0{,}100$~USD/h) jest prawie
dwukrotnie wyższa niż koszt samej pojemności dysku, przez co pod intensywnym obciążeniem
losowym konfiguracja ta jest droższa w~godzinie niż IaaS Premium SSD, a~przy tym $5{,}4$-krotnie
wolniejsza. Tani w~pojemności dysk Standard SSD jest więc nieopłacalny dla obciążeń
transakcyjnych z~dużą liczbą operacji I/O.

\section{Podsumowanie wyników}
\label{sec:wyniki-podsumowanie}

\begin{itemize}
  \item Kolejność wydajności jest jednoznaczna: PaaS General Purpose ($1358$~TPS) $>$ IaaS
  Premium SSD ($936$) $>$ IaaS Standard SSD ($172$) $>$ PaaS Burstable ($121$); dla trzech
  konfiguracji osiągnięto precyzję lepszą niż $\pm5\%$, dla IaaS Standard SSD $\pm8{,}7\%$
  przy sześciu wdrożeniach.
  \item Konfiguracje IaaS są ograniczone przez dysk, konfiguracje PaaS -- przez CPU;
  usługa zarządzana częściowo omija limit dysku dzięki pamięci podręcznej hosta.
  \item Tier dysku decyduje o~wydajności IaaS: Premium SSD P20 daje $5{,}4$-krotnie wyższą
  przepustowość niż Standard SSD E20 o~tej samej pojemności.
  \item Mechanizmy kredytowe zawyżają wydajność w~pierwszych dziesiątkach minut (do
  ok. 33\% dla Premium SSD) lub godzinach (PaaS Burstable) po wdrożeniu; rzetelny pomiar
  stanu utrzymywalnego wymaga ich wyczerpania i~kontroli stanu pul w~każdym przebiegu.
  \item Zmienność między wdrożeniami tej samej konfiguracji przewyższa zmienność między
  powtórzeniami na jednym wdrożeniu, co uzasadnia przyjęcie sesji jako jednostki statystycznej.
  \item Efektywność finansowa PaaS General Purpose i~IaaS Premium SSD jest statystycznie
  nierozróżnialna (ok. 16~mln transakcji na dolara); PaaS Burstable i~IaaS Standard SSD są
  4--7-krotnie mniej efektywne.
\end{itemize}
